% -----------------------------------------------
% Template for ISMIR Papers
% 2026 version, based on previous ISMIR templates

% Requirements :
% * 6+n page length maximum
% * 10MB maximum file size
% * Copyright note must appear in the bottom left corner of first page
% * Clearer statement about citing own work in anonymized submission
% (see conference website for additional details)
% -----------------------------------------------

\documentclass{article}
\usepackage[T1]{fontenc}
\usepackage[utf8]{inputenc}
\usepackage{ismir}
\usepackage{amsmath,amssymb,cite,url}
\usepackage{graphicx}
\usepackage{color}
\graphicspath{{../Draft-legacy/v2-0213-2026/V2-latex/figures/}}

\title{COMPARATIVE ANALYSIS OF PRETRAINED AUDIO REPRESENTATIONS FOR SMALL-SCALE MUSIC ARCHIVES: A CASE STUDY ON THE IPALPITI COLLECTION}

\oneauthor
  {Keita Katsumi}
  {California State Polytechnic University, Pomona\\\texttt{kkatsumi@cpp.edu}}

\def\authorname{K. Katsumi}

\setlength{\parskip}{0pt}
\setlength{\parindent}{1em}

\sloppy % please retain sloppy command for improved formatting

\begin{document}

\maketitle

\begin{abstract}
Learned audio embeddings are increasingly used to support music recommendation and similarity search, yet most existing work and deployed systems are optimized for large, heterogeneous, often pop-centric catalogs.
In contrast, curated classical music archives such as the iPalpiti collection lack large-scale user interaction data and exhibit different structural properties, raising questions about how alternative audio embedding models behave in this domain.
This thesis investigates the performance of multiple backend audio embedding model families on classical music similarity and retrieval tasks, focusing on convolutional, recurrent, hybrid, and transformer-based architectures.
Using a consistent offline evaluation pipeline, raw audio from the iPalpiti archive is transformed into embeddings, indexed for nearest-neighbor search, and evaluated on metadata-defined proxy tasks (e.g., retrieving tracks from the same work or performer).
Ranking quality is assessed with standard information retrieval metrics, including NDCG, Precision@K, and Recall@K.

The experimental study, conducted at modest scale, compares families of pretrained models rather than exhaustively tuning any single architecture.
Results indicate that certain lightweight convolutional and hybrid models can achieve competitive or superior ranking performance to more complex architectures in this classical setting, though differences vary across proxy tasks and evaluation cutoffs.
The analysis highlights architectural and representational choices that most influence ranking quality in this domain, while also illustrating the limitations of offline proxy evaluation.
The thesis concludes with practical recommendations for backend embedding selection in classical music archives and identifies directions for future extension.
\end{abstract}


% Section 1 Introduction
\section{Introduction}\label{sec:introduction}
Modern music streaming platforms rely heavily on machine learning to support search, recommendation,
and discovery.
Much of the prior work in this area, as well as many industrial systems, is built around large-scale
collaborative filtering and deep models trained on extensive user interaction logs.
While effective for mainstream catalogs, these approaches are ill-suited to small, curated archives
that lack dense user data and whose mission emphasizes exploration over pure popularity.
Classical music archives such as the iPalpiti collection exemplify this setting: they contain high-
quality, stylistically coherent recordings but limited usage data, and they require tools that can
surface musically meaningful relationships among works, performances, and artists.

In such environments, content-based recommendation driven by audio embeddings becomes central.
Recent advances in pretrained audio models---including lightweight convolutional networks and large
transformer-based architectures---enable the extraction of rich musical representations from raw
audio.
However, most evaluations of these models are conducted on large-scale, heterogeneous datasets, where
increased model capacity often correlates with improved benchmark performance.
It remains unclear whether this trend transfers to small, single-domain archives such as classical
music collections.
In small-scale settings, limited data diversity, constrained candidate sets, and coarse proxy
definitions may reduce the practical benefit of increased model complexity.
This raises a fundamental design question for content-based recommendation systems in archival
contexts: does increasing embedding-model capacity meaningfully improve similarity retrieval
performance, or does it introduce diminishing returns under constrained conditions?

This thesis addresses this question by focusing exclusively on the backend: the stage that maps audio
recordings to embeddings and supports similarity search, independent of end-to-end user modeling.
Specifically,the study adopts a hypothesis-driven, iterative experimental design.
Rather than assuming that larger models necessarily yield better retrieval performance, we begin with
the hypothesis that increased embedding-model capacity may not improve similarity retrieval in small-
scale, single-domain archives.

To evaluate this hypothesis, retrieval performance is systematically assessed across a spectrum of
model capacities within two representative architectural families (CNN-based and Transformer-based
models).
Proxy retrieval tasks are formulated to approximate musically meaningful similarity, including both a
sanity-check retrieval task (composer-based) and a musical-character-based proxy (e.g., contrasting musical characteristics such as \emph{energetic} vs.\ \emph{calm})
Ranking quality is evaluated using standard information-retrieval metrics (e.g., NDCG, Precision@K, Recall@K).
Experimental conditions are progressively expanded to examine whether higher-capacity models provide
consistent improvements, exhibit diminishing returns, or fail to deliver measurable gains under
constrained archival settings.This thesis is guided by two research questions:

\smallskip
\noindent\textbf{RQ1:} In a small-scale classical music archive (iPalpiti), how does embedding-model
capacity influence content-based music similarity retrieval performance?

\noindent\textbf{RQ2:} Under what conditions, if any, do higher-capacity embedding models provide
meaningful gains over lighter architectures?

\noindent\textbf{Hypothesis:} In a small-scale, single-domain classical music dataset, increasing
embedding-model capacity does not necessarily improve music-similarity retrieval performance.

\smallskip
This work makes three main contributions. First, it provides a systematic empirical investigation of
the relationship between embedding-model capacity and content-based similarity retrieval performance
in a small-scale, single-domain classical archive. Second, it establishes musically grounded
evaluation criteria (via structured proxy retrieval tasks) that enable the identification of
conditions under which retrieval improvements are meaningful. Third, through a hypothesis-driven
experimental design, it analyzes how embedding-model capacity behaves under constrained archival
conditions, clarifying the practical limits of model complexity in small, single-domain
recommendation settings.
\subsection{Thesis Organization}
The remainder of the thesis is organized as follows.
Chapter~2 reviews related work on audio embeddings, evaluation strategies, and system design for
music recommendation.
Chapter~3 formalizes the problem setting and describes the overall backend pipeline.
Chapter~4 details the methodology and experimental design.
Chapter~5 presents and analyzes the results.
Chapter~6 discusses implications, limitations, and directions for future work.













% Section 2 Related Work
\section{Related Work}\label{sec:related-work}

Content-based recommendation is especially important in music collections with sparse or absent user
interaction data, where embeddings derived from the audio signal support similarity search,
retrieval, and downstream recommendation. Recent work has shown the effectiveness of pretrained audio
representations in music information retrieval while also identifying a gap in their systematic
evaluation within recommender systems \cite{tamm2024}. That gap is particularly relevant for small,
domain-specific archives such as classical collections, whose stylistic coherence, limited scale, and
rich structural relationships raise transferability questions for models developed on large,
heterogeneous datasets.

Prior work has explored a broad range of audio representation architectures, including convolutional
neural networks (CNNs), recurrent neural networks (RNNs), transformer-based models, and hybrid
architectures \cite{zaman2023}. Historically, CNN and hybrid CNN–RNN models have dominated supervised
music tasks \cite{jiang2024}, while more recent transformer-based approaches use self-attention to
model longer-range structure \cite{zaman2023}. In this thesis, however, the experimental design
focuses specifically on pure CNN-based and pure transformer-based model families.
CNN-based models remain attractive because their inductive biases support efficient learning of local
spectral structure, and compact CNN architectures have shown strong performance in music detection
and related tasks \cite{reddy2022}. Hybrid architectures have also been explored, including CNN–RNN
models for temporal modeling \cite{lin2024} and CNN–Transformer models for combining local and long-
range representations \cite{pourmoazemi2024}, but they are treated here as background rather than as
primary comparison targets.

Transformer-based models provide a contrasting representation regime: AST applies self-attention
directly to spectrograms \cite{gong2021ast}, and newer large-scale self-supervised systems such as
BEATs rely on pretraining objectives designed to learn general-purpose audio representations from
unlabeled corpora \cite{chen2023beats}. Although these models can deliver strong downstream
performance, their advantages are typically established on large and diverse datasets, leaving open
the question of how well such capacity transfers to small, domain-constrained classical archives.
Evaluation strategy is equally important. Classification benchmarks remain common for audio
embeddings \cite{gong2021ast, lin2024}, but strong classification accuracy does not guarantee well-
structured neighborhoods for similarity retrieval \cite{tamm2024}. For retrieval-oriented settings,
ranking metrics such as NDCG, Precision@K, and Recall@K are more appropriate because they evaluate
the ordering of relevant items directly \cite{krichene2020}. At the same time, these metrics are
sensitive to evaluation design: constrained candidate pools can reduce discriminative power, increase
ties, and even alter comparative conclusions \cite{krichene2020, canamares2020}; this is especially
important in small archival collections. Diversity and long-tail considerations also remain relevant
background concerns in music recommendation, particularly in settings where exploration is valued
\cite{porcaro2022}.

Broader recommender-system work addresses personalization through user modeling, sequential
interaction modeling, and contextual signals \cite{lin2025, abbattista2024, schedl2021}, but those
concerns are outside the scope of this thesis, which isolates backend embedding quality from user
modeling and large-scale deployment constraints. Overall, prior work has established that larger and
more complex models often improve performance in data-rich settings, yet it remains unclear whether
embedding-model capacity provides comparable gains in small, single-domain classical archives. This
thesis addresses that gap by examining how capacity interacts with retrieval-oriented evaluation
under constrained archival conditions \cite{tamm2024, krichene2020, canamares2020}.









% Section 3 Method
\section{Method}\label{sec:method}\label{ch:problem}\label{ch:methodology}
This section formalizes the retrieval setting and experimental framework used to evaluate the effect of embedding-model capacity on similarity-based retrieval performance.

\subsection{Problem Formulation}

Let $D = \{x_1, x_2, \dots, x_N\}$ denote a collection of $N$ classical music recordings drawn from
the iPalpiti archive. Each item $x_i$ represents a complete audio track.

An embedding model is defined as a function:
\[
f_\theta : \mathcal{X} \rightarrow \mathbb{R}^d
\]
where $\theta$ denotes the model parameters and $d$ the embedding dimensionality. The embedding of
item $x_i$ is:
\[
z_i = f_\theta(x_i)
\]

Similarity between two items is computed using a similarity function $s(\cdot,\cdot)$, such as cosine
similarity:
\[
s(z_q, z_i) = \frac{z_q \cdot z_i}{\|z_q\|\|z_i\|}
\]
Higher similarity values indicate greater embedding-space proximity.

Given a query item $x_q \in D$, all items in the collection except $x_q$ are ranked in descending
order of similarity $s(z_q, z_i)$. Retrieval quality is evaluated with respect to a relevance
function $r_q(x_i) \in \{0,1\}$ defined by proxy tasks, and ranking quality is measured using
standard IR metrics including NDCG@K, Precision@K, Recall@K, and F1@K.

The primary independent variable is embedding-model capacity, operationalized mainly through
parameter count ($|\theta|$), with architectural class treated as a separate categorical factor.
Capacity effects are examined within architectural families via small, base, and large variants under
a controlled retrieval setting.

Relevance is defined through two metadata-based proxy tasks. Let $g(x)$ denote a metadata labeling
function.

\textbf{Sanity Proxy.} Relevance is defined by shared composer identity:
\[
r_q(x_i) = 1 \quad \text{if} \quad g_{\text{composer}}(x_q) = g_{\text{composer}}(x_i)
\]

\textbf{Primary Musical Proxy.} Relevance is defined by shared character labels:
\[
r_q(x_i) = 1 \quad \text{if} \quad g_{\text{character}}(x_q) = g_{\text{character}}(x_i)
\]

For experimental tractability, relevance is treated as binary:
\[
r_q(x_i) \in \{0,1\}
\]
with $r_q(x_i)=0$ when labels do not match. This study assumes that such metadata-derived labels
provide a controlled, though coarse, approximation of musical similarity.

\subsection{Retrieval Pipeline}
The backend evaluation pipeline consists of four stages: audio preprocessing, embedding extraction,
similarity computation, and ranking. All components other than embedding extraction are held constant
to isolate the effect of embedding capacity.

The dataset is the iPalpiti classical music archive, comprising [TBD] tracks and [TBD] hours of
audio. Ground truth is derived from editorial metadata, including composer identity and character
labels used in the proxy tasks.

Each complete track serves as the retrieval unit, and embeddings are computed directly from
full-length audio inputs without segmentation. Cosine similarity is used for ranking across all
models.

\subsection{Model Setup}
Models are selected to represent discrete capacity tiers within two architectural families: CNN-based
and Transformer-based models. Within each family, small, medium, and large variants enable controlled
scaling comparisons, while cross-family comparisons are interpreted descriptively because
architecture, parameter count, and pretraining objectives may co-vary.

CNN-based models are drawn from a unified convolutional lineage in which scaling is achieved
primarily through increased depth under comparable training objectives. Specifically, we use Cnn6
(small), Cnn10 (medium), and Cnn14 (large) to represent controlled capacity scaling within the CNN
family.

Transformer-based models are selected from a self-supervised audio representation framework with
documented parameter tiers. We use MERT-95M (medium) and MERT-330M (large) to represent different
capacity levels within the transformer family.

All selected models provide pretrained checkpoints and support standardized embedding extraction.


\subsection{Evaluation Protocol}

Each track in the archive is used as a query in a leave-one-out retrieval setting, where the query
item itself is excluded from the candidate pool. Relevance is determined by metadata matching as
defined above, and performance metrics are averaged across all query instances.

Because relevance labels are binary and the dataset is relatively small, ranking ties may occur when
multiple relevant items share identical similarity scores. To ensure deterministic evaluation, a
consistent tie-breaking strategy (e.g., lexicographic ordering by track identifier) is applied across
all models. Performance is primarily assessed using NDCG@10 and Precision@10, with interpretation
focused on consistent scaling patterns rather than small numerical differences.











% Section 4 Experimental Setup
\section{Experimental Setup}\label{sec:experimental-setup}
This section presents the empirical results of the retrieval experiments defined in Section~\ref{sec:method}. The analysis focuses on the relationship between embedding-model capacity and retrieval performance across two proxy tasks. The primary objective is to determine whether increased model capacity yields monotonic performance gains or exhibits saturation and degradation within the constrained context of a small-scale classical music archive.

\subsection{Experimental Setup Recap}

The experiments were conducted on the full dataset of $N$ tracks from the iPalpiti archive. 
A total of five embedding models were evaluated across two architectural families and three capacity tiers (small, medium, and large).

\begin{itemize}
    \item CNN-Based Family: Convolutional architectures spanning compact to high-capacity variants.
    \item Transformer-Based Family: Transformer architectures spanning medium to large pretrained checkpoints.
\end{itemize}

Performance is primarily evaluated using ranking-based metrics, including Normalized Discounted Cumulative Gain at rank 5 (NDCG@5) and Recall@10. Precision@10 is also reported to reflect the purity of retrieved results at the top ranks. To summarize the trade-off between Precision and Recall, we additionally report F1@10 as the harmonic mean of the two metrics. However, because F1@10 does not account for rank position, it is interpreted as a complementary metric rather than the primary retrieval indicator. 
Given the relatively small dataset size, interpretation focuses on consistent scaling patterns across capacity tiers rather than marginal numerical differences.



% Section 5 Results
\section{Results}\label{sec:results}\label{ch:results}
\subsection{Sanity Proxy Results}

The Sanity Proxy task evaluates the ability of embedding models to retrieve tracks sharing the same composer. This task serves as a baseline for verifying that models capture fundamental compositional structures.

Table~\ref{tab:sanity_results} presents retrieval performance grouped by architectural family and ordered by capacity tier within each family.

\begin{table}[!htbp]
\centering
\caption{Sanity Proxy Results (Composer Retrieval). Metrics reported are corpus-level mean values at rank 10.}
\label{tab:sanity_results}
\resizebox{\columnwidth}{!}{%
\begin{tabular}{lcccc}
\hline
\textbf{Model} & \textbf{NDCG@5} & \textbf{Precision@5} & \textbf{Recall@5} & \textbf{F1@5} \\
\hline
\multicolumn{5}{c}{\textbf{CNN Family}} \\
\hline
CNN-Small & -- & -- & -- & -- \\
CNN-Medium  & -- & -- & -- & -- \\
CNN-Large & -- & -- & -- & -- \\
\hline
\multicolumn{5}{c}{\textbf{Transformer Family}} \\
\hline
Transformer-Medium  & -- & -- & -- & -- \\
Transformer-Large & -- & -- & -- & -- \\
\hline
\end{tabular}
}
\end{table}



The observed trend suggests that [TBD: Description of monotonic/non-monotonic trends within families]. Within the CNN family, performance [increases/decreases/plateaus] as capacity increases. Within the Transformer family, performance [increases/decreases/plateaus] relative to parameter growth.

\subsection{Primary Musical Proxy Results}

The Primary Musical Proxy task evaluates retrieval based on abstract musical character labels (e.g., "Energetic" vs. "Calm"). This task represents the core retrieval challenge for the archive.

Table~\ref{tab:musical_results} summarizes the performance across all models.

\begin{table}[!htbp]
\centering
\caption{Primary Musical Proxy Results (Character Retrieval). Metrics reported are corpus-level mean values at rank 10.}
\label{tab:musical_results}
\resizebox{\columnwidth}{!}{%
\begin{tabular}{lcccc}
\hline
\textbf{Model} & \textbf{NDCG@5} & \textbf{Precision@5} & \textbf{Recall@5} & \textbf{F1@5} \\
\hline
\multicolumn{5}{c}{\textbf{CNN Family}} \\
\hline
CNN-Small & 0.6425 & 0.3895 & 0.3487 & 0.3576 \\
CNN-Medium & 0.6117 & 0.4211 & 0.3575 & 0.3809 \\
CNN-Large & 0.6624 & 0.3579 & 0.3158 & 0.3245\\
\hline
\multicolumn{5}{c}{\textbf{Transformer Family}} \\
\hline
Transformer-Medium  & 0.7087 & 0.4631 & 0.3728 & 0.4059 \\
Transformer-Large & 0.6813 & 0.4526 & 0.3618 & 0.3949 \\
\hline
\end{tabular}
}
\end{table}



In contrast to the Sanity Proxy, trends in the Musical Proxy task suggest [TBD: Describe difference or similarity in trends]. The relationship between capacity and performance appears to be [stronger/weaker] for this more abstract retrieval task. Notably, the highest-capacity models [do/do not] strictly outperform mid-sized architectures.

\subsection{Capacity–Performance Trends Within Architectural Families}
This section analyzes scaling behavior strictly within architectural families to isolate the effect of parameter count from architectural inductive bias.
Taken together, the within-family scaling patterns described below provide the primary empirical test of the condition-dependent capacity hypothesis.


\subsubsection{CNN Scaling Behavior}
Within the CNN family, increasing parameter count from [Model A] to [Model C] results in a [monotonic improvement/saturation/degradation] in NDCG@10. 
Specifically, performance improves from [--] to [--] between the small and base tiers, and from [--] to [--] between the base and large tiers.

\subsubsection{Transformer Scaling Behavior}
Within the Transformer family, scaling from [Model D] to [Model F] exhibits a [monotonic/non-monotonic] trend. While [Model E] achieves a score of [--], the largest model [Model F] achieves [--]. This observation indicates that for transformer-based architectures in this domain, additional capacity [yields diminishing returns/continues to improve performance/degrades performance], consistent with the hypothesis of condition-dependent scaling.

\subsection{Cross-Family Descriptive Comparison}

Comparing performance across families reveals distinct behavioral patterns. The best-performing CNN model ([Model Name]) achieves an NDCG@10 of [--], while the best-performing Transformer ([Model Name]) achieves [--].

It is observed that [Compact CNNs/Large Transformers] tend to perform [better/worse/comparably] in the Sanity Proxy task, whereas in the Musical Proxy task, the performance gap [widens/narrows/reverses]. These comparisons are descriptive and do not imply causal superiority of one architecture over another, as parameter count, architecture, and pretraining data scale co-vary between families. However, the data demonstrates that [Low/High] capacity models from the [CNN/Transformer] family are sufficient to achieve competitive retrieval performance in this specific archival setting.

Given the limited scale of the dataset, minor variations in metric values are expected. Therefore, emphasis is placed on consistent trends across capacity tiers rather than isolated numerical differences.

\subsection{Summary of Empirical Findings}

The experimental results yield the following key observations:

\begin{enumerate}
    \item Within-Family Scaling: Increasing model capacity [does/does not] guarantee monotonic performance improvement. Saturation effects were observed in the [CNN/Transformer] family beyond [--] parameters.
    \item Task Dependence: Capacity scaling behavior [varies/remains consistent] between the Sanity Proxy and the Musical Proxy tasks.
    \item Efficiency: Compact models [competitively match/lag behind] larger models, challenging the assumption that massive parameter counts are strictly necessary for this domain.
\end{enumerate}

Collectively, these findings indicate that embedding-model capacity does not operate as a simple monotonic driver of retrieval quality in small-scale, single-domain archives. Instead, performance patterns vary by task definition and architectural family, supporting a condition-dependent interpretation of scaling behavior.


% Section 6 Discussion
\section{Discussion}\label{sec:discussion}
\subsection{Overview of Findings}
% Briefly restate the purpose of the study and summarize the main findings
% from Chapter 5 at a high level, without repeating numerical results.

\subsection{Implications for Backend Design in Classical Music Archives}
% Discuss what the observed results imply for designing
% embedding-based backends in classical music archives such as iPalpiti.
% Focus on model family selection and practical trade-offs.

\subsection{Limitations}


\subsection{Directions for Future Work}



% Section 7 Conclusion
\section{Conclusion}\label{sec:conclusion}
% Provide a concise concluding statement.
% Revisit Research Question 1 and summarize how it was addressed.
% Emphasize the contribution and potential impact of the work.

% For BibTeX users:
\bibliography{../Draft-legacy/v2-0213-2026/V2-latex/bibliography/references}

\end{document}
